\documentclass[11pt,a4paper]{article}
\usepackage{fontspec,amsmath,amssymb,amsthm,microtype}
\setmainfont{lmroman10-regular.otf}[BoldFont=lmroman10-bold.otf,ItalicFont=lmroman10-italic.otf,BoldItalicFont=lmroman10-bolditalic.otf]
\usepackage[margin=25mm]{geometry}
\usepackage[unicode,hidelinks]{hyperref}
\hypersetup{pdfauthor={LangevinLean Project}}
\setlength{\emergencystretch}{3em}
\setlength{\parskip}{0.4em}
\setlength{\parindent}{0pt}
\allowdisplaybreaks[2]
\raggedbottom
\title{\texorpdfstring{Assumptions and scope of the entropy arguments}{Assumptions and scope of the entropy arguments}}
\author{LangevinLean Project}
\date{}
\begin{document}\maketitle
Draft mathematical exposition; these notes may contain mistakes.


The conditional entropy theorem assumes interpolation densities, their weak equations, and the differentiability and integrability needed for entropy calculus. The Gaussian interpolation theorem constructs these properties under its stated smoothness, logarithmic Sobolev, step, and initial-entropy hypotheses. The assumption check records the different scopes and confirms that finiteness and projection are used with their hypotheses.


\section{\texorpdfstring{Assumed and constructed interpolation regularity}{Assumed and constructed interpolation regularity}}


For the conditional entropy result, smoothness and a lifted logarithmic Sobolev inequality are supplemented by explicit density and weak-equation regularity. The recurrence follows from those data, but that result alone does not prove their existence. The separate Gaussian interpolation construction supplies them for the unconditional theorem. The dependency check retains this distinction, including the finite-divergence requirement before real conversion and the same-map projection requirement.




\end{document}
